\documentclass[tikz,border=3pt]{standalone}

\usepackage[T1]{fontenc}
\usepackage{libertine}
\usepackage{newtxmath}
\usepackage{xcolor}
\usetikzlibrary{arrows.meta,calc,positioning}

\renewcommand{\familydefault}{\sfdefault}

\definecolor{ink}{HTML}{20313F}
\definecolor{muted}{HTML}{667784}
\definecolor{line}{HTML}{A5B3BC}
\definecolor{panelbg}{HTML}{F8FAFB}
\definecolor{blue}{HTML}{2A6F97}
\definecolor{bluefill}{HTML}{EAF3F8}
\definecolor{teal}{HTML}{2A9D8F}
\definecolor{tealfill}{HTML}{E8F6F3}
\definecolor{orange}{HTML}{D97745}
\definecolor{orangefill}{HTML}{FFF1E8}

\tikzset{
  panel/.style={
    rounded corners=5pt,
    draw=line,
    line width=0.75pt,
    fill=panelbg,
    minimum height=5.25cm,
    align=center,
    text=ink,
    inner sep=0pt
  },
  stepbadge/.style={
    circle,
    minimum size=5.0mm,
    inner sep=0pt,
    text=white,
    font=\sffamily\bfseries\scriptsize
  },
  flow/.style={
    -{Latex[length=2.1mm,width=1.45mm]},
    draw=ink,
    line width=0.8pt,
    rounded corners=2pt
  },
  blueflow/.style={
    -{Latex[length=2.1mm,width=1.45mm]},
    draw=blue,
    line width=0.95pt,
    rounded corners=2pt
  },
  orangeflow/.style={
    -{Latex[length=2.1mm,width=1.45mm]},
    draw=orange,
    line width=0.95pt,
    rounded corners=2pt
  },
  chipblue/.style={
    rounded corners=3pt,
    draw=blue,
    fill=bluefill,
    line width=0.7pt,
    text=blue,
    inner xsep=4pt,
    inner ysep=2.5pt,
    font=\sffamily\scriptsize
  },
  chiporange/.style={
    rounded corners=3pt,
    draw=orange,
    fill=orangefill,
    line width=0.7pt,
    text=orange,
    inner xsep=4pt,
    inner ysep=2.5pt,
    font=\sffamily\scriptsize
  },
  railnote/.style={
    fill=white,
    rounded corners=2pt,
    inner xsep=3pt,
    inner ysep=1.5pt,
    font=\sffamily\scriptsize
  }
}

\begin{document}
\begin{tikzpicture}[font=\sffamily, every node/.append style={outer sep=0pt}]
  % Five compact panels mirror the five operations described in the paper.
  \node[panel, minimum width=3.20cm] (prior) at (1.72,2.72) {};
  \node[panel, minimum width=3.20cm] (predict) at (5.22,2.72) {};
  \node[panel, minimum width=3.20cm] (observe) at (8.72,2.72) {};
  \node[panel, minimum width=4.55cm] (update) at (12.90,2.72) {};
  \node[panel, minimum width=3.20cm] (score) at (17.08,2.72) {};

  % Panel titles and numbered badges.
  \node[anchor=north, font=\sffamily\bfseries\scriptsize, text=ink]
    at ($(prior.north)+(0,-0.20)$) {PRIOR STATE};
  \node[anchor=north, font=\sffamily\bfseries\scriptsize, text=ink]
    at ($(predict.north)+(0,-0.20)$) {PREDICT};
  \node[anchor=north, font=\sffamily\bfseries\scriptsize, text=ink]
    at ($(observe.north)+(0,-0.20)$) {LABEL RESPONSE};
  \node[anchor=north, font=\sffamily\bfseries\scriptsize, text=ink]
    at ($(update.north)+(0,-0.20)$) {BAYESIAN UPDATE};
  \node[anchor=north, font=\sffamily\bfseries\scriptsize, text=ink]
    at ($(score.north)+(0,-0.20)$) {STYLE SCORE};

  \node[stepbadge, fill=blue] at ($(prior.north west)+(0.17,-0.17)$) {1};
  \node[stepbadge, fill=blue] at ($(predict.north west)+(0.17,-0.17)$) {2};
  \node[stepbadge, fill=orange] at ($(observe.north west)+(0.17,-0.17)$) {3};
  \node[stepbadge, fill=teal] at ($(update.north west)+(0.17,-0.17)$) {4};
  \node[stepbadge, fill=teal] at ($(score.north west)+(0.17,-0.17)$) {5};

  % Step 1: prior state distribution b_{t-1}.
  \node[font=\scriptsize, text=muted] at (1.72,4.58)
    {state belief \(b_{t-1}(s)\)};
  \draw[draw=line, line width=0.6pt] (1.62,1.10) -- (1.62,4.16);

  \node[anchor=east, align=right, font=\scriptsize] at (1.49,3.78)
    {Emotion\\disclosure};
  \fill[blue] (1.66,3.66) rectangle ++(0.66,0.24);
  \node[anchor=west, font=\scriptsize] at (2.42,3.78) {0.30};

  \node[anchor=east, align=right, font=\scriptsize] at (1.49,2.82)
    {Emotion\\processing};
  \fill[blue] (1.66,2.70) rectangle ++(0.99,0.24);
  \node[anchor=west, font=\scriptsize] at (2.75,2.82) {0.45};

  \node[anchor=east, align=right, font=\scriptsize] at (1.49,1.86)
    {Considering\\solutions};
  \fill[blue] (1.66,1.74) rectangle ++(0.55,0.24);
  \node[anchor=west, font=\scriptsize] at (2.31,1.86) {0.25};

  % Step 2: transition-model prediction.
  \node[chipblue] at (5.22,4.42)
    {transition model \(P(s_t\mid s_{t-1})\)};
  \node[
    circle,
    draw=blue,
    fill=white,
    minimum size=5.5mm,
    font=\scriptsize
  ] (s1) at (4.35,3.55) {\(s_1\)};
  \node[
    circle,
    draw=blue,
    fill=white,
    minimum size=5.5mm,
    font=\scriptsize
  ] (s2) at (5.22,3.55) {\(s_2\)};
  \node[
    circle,
    draw=blue,
    fill=white,
    minimum size=5.5mm,
    font=\scriptsize
  ] (s3) at (6.09,3.55) {\(s_3\)};
  \draw[-{Latex[length=1.5mm]}, draw=blue, line width=0.65pt] (s1) -- (s2);
  \draw[-{Latex[length=1.5mm]}, draw=blue, line width=0.65pt] (s2) -- (s3);
  \draw[-{Latex[length=1.5mm]}, draw=blue, line width=0.55pt, bend left=42]
    (s3.north west) to (s1.north east);

  \node[align=center, font=\scriptsize] at (5.22,2.63) {%
    \(\displaystyle
      \begin{aligned}
        \hat b_t(s_t)
        &=\sum_{s_{t-1}\in S}P(s_t\mid s_{t-1})\\[-1pt]
        &\hspace{7mm}{}\cdot b_{t-1}(s_{t-1})
      \end{aligned}
    \)
  };
  \node[chipblue] (predout) at (5.22,1.28)
    {predicted distribution \(\hat b_t(s_t)\)};

  % Step 3: candidate response is mapped to an observation label.
  \node[
    rounded corners=3pt,
    draw=orange,
    fill=white,
    line width=0.7pt,
    text width=2.42cm,
    align=center,
    inner xsep=4pt,
    inner ysep=5pt,
    font=\scriptsize
  ] (response) at (8.72,4.10)
    {\textit{``That sounds really difficult.''}};
  \node[font=\scriptsize, text=muted] at (8.72,3.45)
    {candidate response \(r_t\)};
  \node[
    circle,
    draw=orange,
    fill=orangefill,
    minimum size=8.5mm,
    font=\sffamily\bfseries\scriptsize,
    text=orange
  ] (llm) at (8.72,2.72) {LLM};
  \draw[orangeflow] (response.south) -- (llm.north);
  \node[chiporange, align=center] (obsout) at (8.72,1.35)
    {observed label \(o_t\)\\[-1pt]\textbf{Empathy / validation}};
  \draw[orangeflow] (llm.south) -- (obsout.north);

  % Step 4: the two variables enter separate terms in the Bayesian update.
  \node[chiporange] (obsport) at (11.92,4.55)
    {observation \(o_t\)};
  \node[chipblue] (predport) at (14.00,4.55)
    {prediction \(\hat b_t\)};

  \node[font=\small] at (11.32,3.68) {\(b_t(s_t)=\)};
  \node[font=\small, text=orange] (liketerm) at (12.55,3.76)
    {\(P(o_t\mid s_t)\)};
  \node[font=\small, text=muted] at (13.25,3.76) {\(\times\)};
  \node[font=\small, text=blue] (predterm) at (14.03,3.76)
    {\(\hat b_t(s_t)\)};
  \draw[draw=ink, line width=0.55pt] (12.08,3.46) -- (14.82,3.46);
  \node[font=\scriptsize] at (13.45,3.08) {%
    \(\displaystyle
      \sum_{s'\in S}P(o_t\mid s')\,\hat b_t(s')
    \)
  };
  \draw[orangeflow] (obsport.south) -- (liketerm.north);
  \draw[blueflow] (predport.south) -- (predterm.north);

  \node[font=\scriptsize, text=muted] at (12.90,2.55)
    {posterior state belief \(b_t(s)\)};
  \draw[draw=line, line width=0.55pt] (13.24,0.64) -- (13.24,2.29);

  \node[anchor=east, font=\scriptsize] at (13.10,2.05) {Disclosure};
  \fill[muted] (13.28,1.95) rectangle ++(0.26,0.20);
  \node[anchor=west, font=\scriptsize] at (13.62,2.05) {0.12};

  \node[anchor=east, font=\scriptsize, text=teal] at (13.10,1.48) {Processing};
  \fill[teal] (13.28,1.38) rectangle ++(1.45,0.20);
  \node[anchor=west, font=\scriptsize] at (14.81,1.48) {0.66};

  \node[anchor=east, font=\scriptsize, text=teal] at (13.10,0.91) {Solutions};
  \fill[teal] (13.28,0.81) rectangle ++(0.48,0.20);
  \node[anchor=west, font=\scriptsize] at (13.84,0.91) {0.22};

  % Explicit input rails preserve the informative arrows from the source slide
  % while keeping them outside every text and equation block.
  \draw[blueflow]
    (predout.east)
    -- (6.93,1.28)
    -- (6.93,6.08)
    -- node[railnote, text=blue, pos=0.66]
      {predicted state distribution}
    (14.00,6.08)
    -- (predport.north);
  \draw[orangeflow]
    (obsout.east)
    -- (10.48,1.35)
    -- (10.48,5.78)
    -- node[railnote, text=orange, pos=0.50]
      {response-derived observation}
    (11.92,5.78)
    -- (obsport.north);

  % Step 5: sum only the target-state posterior mass.
  \node[align=center, font=\scriptsize] at (17.08,4.27) {%
    \(\displaystyle
      \mathrm{Score}(r_t)=
      \sum_{s\in S_{\mathrm{target}}} b_t(s)
    \)
  };
  \node[
    rounded corners=3pt,
    draw=teal,
    fill=tealfill,
    text=teal,
    text width=2.35cm,
    align=center,
    inner xsep=4pt,
    inner ysep=4pt,
    font=\scriptsize
  ] at (17.08,3.22) {%
    \(S_{\mathrm{target}}\)\\
    Processing \(+\) Solutions
  };
  \node[font=\small, text=ink] at (17.08,2.20)
    {\(0.66+0.22\)};
  \node[font=\bfseries\Large, text=teal] at (17.08,1.34)
    {\(\mathbf{0.88}\)};
  \node[font=\scriptsize, text=muted] at (17.08,0.78)
    {style compatibility};

  % Primary state flow: prior -> prediction, posterior -> score.
  \draw[flow] (prior.east) -- (predict.west);
  \draw[flow] (update.east) -- node[
    fill=white,
    text=muted,
    font=\scriptsize,
    inner xsep=2pt,
    inner ysep=1pt
  ] {posterior} (score.west);
\end{tikzpicture}
\end{document}
